The closecapture instruction pops an action from the action list,
and pushes a capture on the capture list.
The popped action must correspond to the given register value,
and the pushed capture will be filled with the combined data from the
current state, and the popped action.

\begin{myquote}
\begin{verbatim}

when state = S_EXECUTE =>
  if opcode = OP_CLOSECAPTURE then
    if actioncount = 0 then
      errorcode <= ERR_ACTIONCOUNT;
      state <= S_ERROR;
      v_redirected := true; -- prevent fetchop
    else
      actionpopped <= actionmem(to_integer(actioncount - 1));
      if actionpopped.reg /= instrreg then
        errorcode <= ERR_ACTIONMISMATCH;
        state <= S_ERROR;
        v_redirected := true; -- prevent fetchop
      else
        pushcapturetuple <= (
          reg      = actionpopped.reg,
          offset   = actionpopped.offset,
          size     = inputptr - actionpopped.offset,
          stacklen = stackcount
        );
        actioncount <= actioncount - 1;
        state <= S_PUSHCAPTURE;
        v_redirected := true; -- prevent fetchop
      end if;
    end if;
  end if;

\end{verbatim}
\end{myquote}

\begin{myquote}
\begin{verbatim}

when state == S_PUSHCAPTURE:
  if capturecount < max then
    capturemem[ capturecount ] <= pushcapturetuple;
    capturecount <= capturecount + 1;
    state <= S_FETCHOP; // end of instruction handling
  else
    errorcode <= ERR_CAPTURECOUNT;
    state <= S_ERROR;
  endif;

\end{verbatim}
\end{myquote}
